\originalchapter{条件概率、Bayes 公式与独立性}
\sourcelecture{6--7}

一个试验可能已经完成, 但我们只得到部分信息. 掷骰之后有人告诉我们结果是偶数, 样本空间中的奇数结果就可以排除. 条件概率把这次信息更新写成一个新的概率模型. 它也让我们能够按时间顺序分解复杂试验, 或者把``某个原因产生某种观察''的概率反过来用于推断原因.

\originalsection{信息限制与归一化}
\begin{definition}
若 $F\in\mathcal F$ 且 $\Prob(F)>0$, 定义事件 $E$ 在给定 $F$ 下的条件概率为
\[
 \Prob(E\mid F)=\frac{\Prob(E\cap F)}{\Prob(F)}.
\]
竖线右边是已知发生的事件, 左边是仍要判断的事件.
\end{definition}
在有限均匀模型中, 这个比值是 $|E\cap F|/|F|$, 相当于把 $F$ 当成新的样本空间. 在非均匀模型中, 排除 $F$ 外部的结果之后, 内部结果的相对权重仍然保持, 统一除以 $\Prob(F)$ 使总概率重新成为一. 因而信息更新并不自动把剩余结果变成等可能.

\begin{proposition}
固定正概率事件 $F$, 映射 $E\mapsto\Prob(E\mid F)$ 是 $(\Omega,\mathcal F)$ 上的概率测度. 若把样本空间限制为 $F$, 则它也在事件族 $\{E\cap F:E\in\mathcal F\}$ 上定义一个概率空间.
\end{proposition}
\begin{proof}
单调性给 $0\leq\Prob(E\cap F)\leq\Prob(F)$, 所以条件概率在 $[0,1]$ 内. 又 $\Prob(\Omega\mid F)=\Prob(F)/\Prob(F)=1$. 对两两互斥的 $E_n$, 各 $E_n\cap F$ 也互斥, 且
\[
 \Prob\left(\bigcup_n E_n\mid F\right)
 =\frac{\Prob(\bigcup_n(E_n\cap F))}{\Prob(F)}
 =\sum_n\frac{\Prob(E_n\cap F)}{\Prob(F)}.
\]
所以可数可加性成立. 限制到 $F$ 后, 补集相对于 $F$ 取, 并集仍留在事件族内; 如果两个事件与 $F$ 的交相同, 它们给出相同条件概率, 因而定义无歧义.
\end{proof}
因此前一章的补集、并集、容斥公式都可以在条件概率模型中使用. 例如 $\Prob(E^c\mid F)=1-\Prob(E\mid F)$. 同时, $\Prob(F\mid F)=1$ 并不表示 $F=\Omega$, 只表示新模型集中在 $F$ 上.

\begin{example}
公平十面骰的点数为 $1,\ldots,10$. 已知点数大于四, 求它大于七的概率. 再求已知大于七时大于四的概率.
\end{example}
\begin{solution}
令 $E=\{8,9,10\}$, $F=\{5,6,7,8,9,10\}$. 有 $E\subseteq F$, 所以 $\Prob(E\mid F)=(3/10)/(6/10)=1/2$. 反过来 $\Prob(F\mid E)=1$. 条件概率一般不对称; 它的分母由已知的信息决定.
\end{solution}

当 $\Prob(F)=0$ 时, 上面的定义是 $0/0$, 没有定义. 不能直接宣布它为零或一. 以后对连续随机变量的条件分布, 会借助密度或条件期望建立另一种定义, 其适用条件和版本需要单独说明. 本章所有使用比值的地方都要求分母为正.

\originalsection{乘法公式与多步试验}
把条件概率定义重新排列可得
\[
 \Prob(E\cap F)=\Prob(F)\Prob(E\mid F).
\]
这条式子不要求 $E,F$ 独立. 它的意义是先达到 $F$, 再在 $F$ 已发生的模型里达到 $E$.

\begin{theorem}[链式乘法公式]
若 $E_1,\ldots,E_n$ 的每个前缀交集 $E_1\cap\cdots\cap E_j$ 在 $1\leq j<n$ 时具有正概率, 则
\[
 \Prob\left(\bigcap_{j=1}^n E_j\right)
 =\Prob(E_1)\prod_{j=2}^n
 \Prob\left(E_j\mid\bigcap_{i=1}^{j-1}E_i\right).
\]
\end{theorem}
\begin{proof}
记 $F_j=\bigcap_{i=1}^jE_i$. 定义给 $\Prob(F_j)=\Prob(F_{j-1})\Prob(E_j\mid F_{j-1})$. 从 $j=2$ 开始逐次代入, 中间的 $\Prob(F_j)$ 消去, 得所述乘积. 若某个前缀的概率为零, 完整交集也因包含于它而概率为零, 但后面的条件概率可能没有定义; 此时应在零概率前缀处停止计算, 不把未定义的项写入乘积.
\end{proof}

\begin{example}
袋中有五个红球、三个蓝球, 不放回依次取三球. 求依次出现红、蓝、红的概率, 并与恰有两红一蓝的概率比较.
\end{example}
\begin{solution}
链式公式给
\[
 \Prob(RBR)=\frac58\frac37\frac46=\frac5{28}.
\]
第二项分母为七、第三项分母为六, 因为前面的球已被取走; 第三项红球数为四, 因为第一步已取出一个红球. 另两种次序 $RRB,BRR$ 有相同概率, 所以恰有两红一蓝的概率为 $15/28$. 也可在无序空间验证为 $\binom52\binom31/\binom83=15/28$. 把每一步红球概率都写成 $5/8$ 则误用了独立性.
\end{solution}

\begin{example}
均匀排列 $n$ 个元素. 已知指定的前 $r$ 个元素都固定, 求另一个指定元素也固定的概率, 其中 $0\leq r<n$.
\end{example}
\begin{solution}
已有固定条件的排列数为 $(n-r)!$, 增加一个固定条件后为 $(n-r-1)!$. 两者之比为 $1/(n-r)$. 因而指定 $k$ 个元素同时固定的概率为
\[
 \frac1n\frac1{n-1}\cdots\frac1{n-k+1}=\frac{(n-k)!}{n!}.
\]
这重新得到错排计算中的交集概率, 并说明固定事件的条件概率会随已有固定数上升.
\end{solution}

\originalsection{全概率公式与 Bayes 公式}
当一个观察可能由几种情形产生时, 应分别数出各情形的贡献, 而不是把各条件概率不加权地平均.

\begin{theorem}[全概率公式]
设 $H_1,H_2,\ldots$ 是有限或可数个两两互斥的事件, 其并为 $\Omega$, 且各 $\Prob(H_i)>0$. 则对任意事件 $B$,
\[
 \Prob(B)=\sum_i\Prob(B\mid H_i)\Prob(H_i).
\]
若 $\Prob(B)>0$, 则
\[
 \Prob(H_j\mid B)=
 \frac{\Prob(B\mid H_j)\Prob(H_j)}
 {\sum_i\Prob(B\mid H_i)\Prob(H_i)}.
\]
\end{theorem}
\begin{proof}
$B=\bigcup_i(B\cap H_i)$ 是互斥分解. 可数可加性和两事件乘法公式给出第一式. 条件概率定义给
\[
 \Prob(H_j\mid B)=\frac{\Prob(H_j\cap B)}{\Prob(B)}
 =\frac{\Prob(B\mid H_j)\Prob(H_j)}{\Prob(B)},
\]
再用第一式展开分母便得第二式. 若划分中有零概率事件, 它们的交集贡献为零, 可以先删去, 无须对它们定义条件概率.
\end{proof}
第二式叫 Bayes 公式. $\Prob(H_j)$ 是观察之前的先验概率; $\Prob(B\mid H_j)$ 是情形 $H_j$ 下产生观察 $B$ 的似然; $\Prob(H_j\mid B)$ 是观察之后的后验概率. Bayes 公式没有额外创造一个关于世界的假设, 它从概率公理推出. 真正需要建模的是先验和似然, 以及观测事件到底是什么.

\begin{example}
一种自动质检系统用于识别瑕疵品. 产品中瑕疵率为 $1\%$. 对瑕疵品, 系统报警概率为 $0.95$; 对合格品, 系统误报概率为 $0.04$. 已知系统报警, 求产品确有瑕疵的概率.
\end{example}
\begin{solution}
令 $D$ 为瑕疵事件, $T$ 为报警. 划分 $D,D^c$ 给
\[
 \Prob(T)=0.95\cdot0.01+0.04\cdot0.99=0.0491.
\]
因而
\[
 \Prob(D\mid T)=\frac{0.0095}{0.0491}=\frac{95}{491}\approx0.1935.
\]
可以想象十万件产品: 约一千件有瑕疵, 其中九百五十件报警; 九万九千件合格, 其中三千九百六十件误报. 报警产品中只有九百五十件有瑕疵. 这解释了为什么较高的识别率也可能伴随较低的报警后瑕疵率. ``准确率''一词若未区分这两个条件概率及误报概率, 就不足以解题.
\end{solution}

\begin{example}
两条生产线供给同一种零件. 甲线占总供给的 $1/5$, 乙线占 $4/5$. 检测人员将甲线零件正确识别为甲的概率为 $0.9$, 将乙线零件误识别为甲的概率为 $0.1$. 给定检测报告为甲, 求真实来自甲线的概率.
\end{example}
\begin{solution}
真实来源必须按实际供给占比加权. 所求为
\[
 \frac{0.9/5}{0.9/5+0.1\cdot4/5}=\frac9{13}.
\]
报告的正确识别率 $0.9$ 是给定真实来源后的概率, 不能直接作为给定报告后的概率. 若零件进入检测流程的机会还依赖生产线, 应先以进入流程后的供给占比作为先验; 总产量的占比就不一定适用.
\end{solution}

从两事件 Bayes 公式还可看出, 当 $\Prob(A)>0$ 且 $\Prob(B)>0$ 时,
\[
 \frac{\Prob(A\mid B)}{\Prob(A)}=\frac{\Prob(B\mid A)}{\Prob(B)}.
\]
右边等于一表示观察没有改变 $A$ 的概率; 等于零表示观察与 $A$ 不相容, 后验为零; 它最大可到 $1/\Prob(A)$, 这时后验为一. 这些说法都依赖模型中的似然, 并不单凭观察看起来特别就成立.

\originalsection{胜算、连续更新与模型检查}
对 $0<\Prob(A)<1$, 定义 $A$ 的胜算为 $\Prob(A)/\Prob(A^c)$. 它是发生与不发生的概率之比, 不是概率本身. 胜算为 $r$ 时, 概率为 $r/(1+r)$.

\begin{proposition}[胜算更新]
若 $0<\Prob(A)<1$, $\Prob(B\mid A)>0$, $\Prob(B\mid A^c)>0$, 则
\[
 \frac{\Prob(A\mid B)}{\Prob(A^c\mid B)}
 =\frac{\Prob(A)}{\Prob(A^c)}
 \frac{\Prob(B\mid A)}{\Prob(B\mid A^c)}.
\]
\end{proposition}
\begin{proof}
分别对 $A,A^c$ 使用 Bayes 公式. 两个后验概率的分母同为 $\Prob(B)$, 相除后消去, 剩下所述乘积. 似然比的分母为正保证胜算比有意义. 若某个似然为零, 应直接用 Bayes 公式判断后验是否为零或一.
\end{proof}
例如先验胜算为 $2:3$, 观察的似然比为四, 则后验胜算为 $8:3$, 后验概率为 $8/11$. 有多项证据 $B_1,\ldots,B_m$ 时, 正确的似然比是
\[
 \frac{\Prob(B_1\cap\cdots\cap B_m\mid A)}
 {\Prob(B_1\cap\cdots\cap B_m\mid A^c)}.
\]
可以用链式公式将它写为一串\emph{有条件}的似然比. 只有证据在给定 $A$ 后相互独立, 且在给定 $A^c$ 后也相互独立时, 才能改成单项似然比的乘积.

\begin{example}
延续瑕疵品模型, 假设两次检测在给定产品质量后独立, 并具有相同报警率. 求两次都报警之后的瑕疵概率.
\end{example}
\begin{solution}
条件独立给 $\Prob(T_1T_2\mid D)=0.95^2$, $\Prob(T_1T_2\mid D^c)=0.04^2$. 因而
\[
 \Prob(D\mid T_1T_2)=
 \frac{0.01\cdot0.95^2}{0.01\cdot0.95^2+0.99\cdot0.04^2}
 =\frac{0.009025}{0.010609}\approx0.8507.
\]
假如第二次只是重复读取同一个传感器保存的同一结果, 那么 $T_2=T_1$, 两次报警事件等于一次报警事件, 后验仍为 $95/491$. 所以增加报告数量不必然增加独立证据.
\end{solution}
主观先验也可以满足公理并按 Bayes 更新, 但公式不能替输入数字提供事实依据. 还应核对证据的选择过程: 如果只挑选最支持某结论的观察来报告, ``收到这份报告''与``预先指定的观察发生''不是同一事件. 这类偏差应写进观测机制, 而不能靠事后套用公式消除.

\originalsection{主持人问题与信息产生方式}
条件概率需要指定信息怎样产生. 对下面的试验, 开门规则是模型的一部分, 并不是无关的故事细节.

\begin{example}
四个门中恰有一个藏有奖品, 位置均匀. 参加者先选门一. 知道位置的主持人从另外三个门中打开两个无奖品的门; 如果有多种可开的二门集合, 他均匀选择其中一个. 看见主持人打开门二和门三后, 坚持门一与换到门四, 获奖概率各是多少?
\end{example}
\begin{solution}
令 $J_i$ 表示奖品在门 $i$, $O$ 表示打开门二、三. 当奖品在门一时, 三个其余门都无奖品, 有三种可开的二门集合, 所以 $\Prob(O\mid J_1)=1/3$. 当奖品在门四时, 主持人只能开门二、三, 所以 $\Prob(O\mid J_4)=1$. 当奖品在门二或三时, $O$ 不可能发生. 因而
\[
 \Prob(O)=\frac14\frac13+\frac14=\frac13,
 \quad \Prob(J_1\mid O)=\frac14,
 \quad \Prob(J_4\mid O)=\frac34.
\]
换门更有利. 在主持人总会留下恰好一个其他门的规则下, 换门赢恰好对应初选错误, 也直接给 $3/4$. 主持人只公布了自己保证无奖品的门, 这并不是从所有门中无条件删除两个门.

若主持人不知道奖品位置, 只是从门二、三、四中均匀随机打开两个, 则他可能打开有奖品的门. 此时给定他碰巧打开门二、三且两门都空, $J_1$ 和 $J_4$ 对这个观察的似然都为 $1/3$, 后验各为 $1/2$. 两种试验留下相同的画面, 但信息产生机制不同, 条件概率也不同.
\end{solution}

\begin{example}
一个家庭恰有两个孩子, 区分长幼. 为说明计数假设, 把每个孩子的性别建模为独立的男或女, 各概率 $1/2$; 出生月份独立均匀为十二个月, 且与性别独立. 已知家庭中至少有一个男孩出生于一月, 求两个孩子都是男孩的概率. 若改成均匀选一个孩子, 并观察该孩子是生于一月的男孩, 答案是否相同?
\end{example}
\begin{solution}
每个孩子有二十四种等可能的性别月份组合, 有序家庭共有 $24^2=576$ 种. 令 $s$ 为``一月男孩''这个指定类别. 至少一个孩子为 $s$ 的家庭有 $576-23^2=47$ 种. 两个孩子都为男孩且至少一个为 $s$ 的家庭, 在 $12^2$ 个男孩月份对中排除 $11^2$ 个无一月对, 留下 $23$ 种. 所求为 $23/47$.

若均匀选取一个孩子后观察到其类别为 $s$, 还要把``选了长子还是幼子''纳入样本空间. 被选孩子的结果不限制另一个孩子, 另一个仍有一半概率为男孩, 所以答案为 $1/2$. 第一种机制中两个孩子都是 $s$ 的家庭只计一次; 第二种机制中它既可通过选长子报告, 又可通过选幼子报告, 权重不同. 实际家庭数据未必满足均匀月份或上述性别假设; 这里只建立一个明确可算的模型.
\end{solution}

\originalsection{独立事件与互斥事件}
\begin{definition}
事件 $A,B$ 独立, 是指
\[
 \Prob(A\cap B)=\Prob(A)\Prob(B).
\]
当 $\Prob(B)>0$ 时, 这等价于 $\Prob(A\mid B)=\Prob(A)$; 当 $\Prob(A)>0$ 时, 也等价于 $\Prob(B\mid A)=\Prob(B)$.
\end{definition}
等价关系只须代入条件概率定义并乘除正分母. 乘积定义还能处理零概率事件, 所以比条件概率说法更普遍. 独立是一种关于概率的对称关系, 不是关于物理因果的断言. 某事件用到同一个硬币结果, 仍可能与另一个事件独立; 两个看似无关的故事, 也可能在所建模型中不独立.

\begin{proposition}
若 $A,B$ 独立, 则 $A^c,B$, $A,B^c$, $A^c,B^c$ 也独立. 零概率事件及概率一事件与任意事件独立. 若 $A,B$ 互斥且两者概率都为正, 则它们不独立.
\end{proposition}
\begin{proof}
分解 $B=(A\cap B)\cup(A^c\cap B)$ 得
\[
 \Prob(A^c\cap B)=\Prob(B)-\Prob(A\cap B)
 =(1-\Prob(A))\Prob(B)=\Prob(A^c)\Prob(B).
\]
交换 $A,B$ 并再取补可得另外两式. 若 $\Prob(A)=0$, 则单调性给 $\Prob(A\cap B)=0$, 符合独立定义. 概率一事件的结论从其零概率补集得到. 最后, 互斥时交集概率为零, 而正概率的乘积大于零, 因而不能独立.
\end{proof}
互斥描述不能同时发生, 独立描述知道其中一个发生后不改变另一个的概率. 这两个概念在正概率情形下恰好不相容.

\begin{example}
公平抛两枚硬币. 令 $A$ 为第一枚正面, $B$ 为第二枚正面, $C$ 为两枚结果不同. 验证任意两事件独立, 并判断三事件是否相互独立.
\end{example}
\begin{solution}
空间是四个等概率结果 $HH,HT,TH,TT$. 三事件概率都为 $1/2$. $A\cap B=\{HH\}$, $A\cap C=\{HT\}$, $B\cap C=\{TH\}$, 概率都为 $1/4$, 等于各边缘概率乘积. 但是 $A\cap B\cap C=\varnothing$, 概率为零, 与 $(1/2)^3=1/8$ 不同. 知道第一枚正面时, 两枚不同的概率仍为一半; 知道两枚都正面时, 两枚不同的概率却为零.
\end{solution}

\originalsection{成对独立与相互独立}
\begin{definition}
有限事件族 $A_1,\ldots,A_n$ 称为相互独立, 若对每个非空指标子集 $I\subseteq\{1,\ldots,n\}$ 都有
\[
 \Prob\left(\bigcap_{i\in I}A_i\right)=\prod_{i\in I}\Prob(A_i).
\]
若只要求 $|I|=2$ 的这些等式, 则称成对独立. 无限事件族相互独立, 是指它的每个有限子族都相互独立.
\end{definition}
三个事件相互独立时, 要检查三个二重交集和一个三重交集, 单事件条件本身自动成立. 只检查全部事件的交集也不足以保证独立. 上例说明成对独立不推出相互独立, 因为联合信息可以揭示单条信息没有揭示的约束.

\begin{proposition}
相互独立事件族中, 将任意若干事件替换为其补集, 所得事件族仍相互独立. 特别地, 对每个 $\eps_i\in\{0,1\}$, 令 $A_i^{(1)}=A_i$, $A_i^{(0)}=A_i^c$, 则
\[
 \Prob\left(\bigcap_{i=1}^n A_i^{(\eps_i)}\right)
 =\prod_{i=1}^n\Prob(A_i^{(\eps_i)}).
\]
\end{proposition}
\begin{proof}
先只替换 $A_j$. 对任何含 $j$ 的指标子集 $I$, 令 $G=\bigcap_{i\in I\setminus\{j\}}A_i$. 独立条件和互斥分解给
\[
 \Prob(G\cap A_j^c)=\Prob(G)-\Prob(G\cap A_j)
 =\left(\prod_{i\in I\setminus\{j\}}\Prob(A_i)\right)(1-\Prob(A_j)).
\]
这恰为新事件的概率乘积. 不含 $j$ 的子集条件未改变, 所以替换一个后仍相互独立. 逐个替换即可得到一般结论; 空交集 $G=\Omega$ 的概率与空乘积都取一.
\end{proof}
这个命题使独立试验的``成功或失败''序列具有乘积概率. 若 $A_i$ 都有概率 $p$, 则任何指定含 $k$ 个成功的模式概率为 $p^k(1-p)^{n-k}$, 而这类模式有 $\binom nk$ 个, 因而恰好成功 $k$ 次的概率为 $\binom nkp^k(1-p)^{n-k}$. 后面会把它作为二项分布研究.

若 $I,J$ 不相交且 $\Prob(\bigcap_{j\in J}A_j)>0$, 相互独立还给
\[
 \Prob\left(\bigcap_{i\in I}A_i\,\middle|\,\bigcap_{j\in J}A_j\right)
 =\frac{\prod_{i\in I\cup J}\Prob(A_i)}{\prod_{j\in J}\Prob(A_j)}
 =\prod_{i\in I}\Prob(A_i).
\]
由此可见, 相互独立控制的不只是两两信息, 也控制任意一组事件联合发生后的信息更新.

\originalsection{随机次序中的独立与相关}
\begin{example}
将若干不同标签均匀随机排列. 设 $E_{a,b}$ 表示标签 $a$ 出现在 $b$ 之前. 当 $a,b,c,d$ 互不相同时, 验证 $E_{a,b}$ 与 $E_{c,d}$ 独立. 再计算 $\Prob(E_{a,b}\mid E_{a,c})$.
\end{example}
\begin{solution}
交换标签 $a,b$ 是排列空间的双射, 它翻转 $E_{a,b}$ 而保持 $E_{c,d}$; 交换 $c,d$ 同样翻转另一事件. 所以这两个事件的四种真假组合有相同排列数, 各概率为 $1/4$. 每个事件概率为 $1/2$, 因而独立.

对第二问, 只需考察 $a,b,c$ 的相对次序, 六种次序等概率. $E_{a,c}$ 对应 $abc,acb,bac$ 三种, 其中前两种同时满足 $E_{a,b}$, 因而条件概率为 $2/3$, 不等于 $1/2$. 共享一个标签的先后关系在这里不独立.
\end{solution}

\begin{exercise}
在均匀随机排列中, 已知 $a$ 出现在指定的 $r$ 个不同标签之前. 另取一个不在这些标签中且不同于 $a$ 的标签 $b$. 求 $a$ 也出现在 $b$ 之前的条件概率.
\end{exercise}
\begin{solution}
只看这 $r+2$ 个标签的相对次序即可. 每种相对次序等概率: 每个指定相对次序在全部排列中对应相同数目的位置选择和其余标签排列. 已知事件表示 $a$ 是 $r+1$ 个指定标签中最早的, 所以其概率为 $1/(r+1)$. 它与 $a$ 早于 $b$ 的交集表示 $a$ 是全部 $r+2$ 个标签中最早的, 概率为 $1/(r+2)$. 两者相除, 所求为 $(r+1)/(r+2)$. 增加已有的领先信息会提高 $a$ 领先新标签的条件概率.
\end{solution}

\originalsection{练习与解答}
\begin{exercise}
盒子甲有三白一黑四球, 盒子乙有一白三黑四球. 先以概率 $1/3$ 选甲、$2/3$ 选乙, 再不放回从选定盒子取两球. 求两球都白的概率, 以及两球都白后选了甲的概率.
\end{exercise}
\begin{solution}
给定甲时两白概率为 $(3/4)(2/3)=1/2$; 给定乙时两白不可能. 全概率给 $\Prob(WW)=(1/3)(1/2)+(2/3)0=1/6$. Bayes 给 $\Prob(\text{甲}\mid WW)=1$. 一次白球的观察却不能确定盒子: 其甲后验为 $((1/3)(3/4))/((1/3)(3/4)+(2/3)(1/4))=3/5$.
\end{solution}

\begin{exercise}
设 $A,B$ 独立且各有概率 $1/2$. 在已知 $A\cup B$ 发生的条件概率空间中, $A,B$ 是否仍独立?
\end{exercise}
\begin{solution}
原模型 $\Prob(A\cap B)=1/4$, $\Prob(A\cup B)=3/4$. 条件模型中 $\Prob(A\mid A\cup B)=\Prob(B\mid A\cup B)=2/3$, 而 $\Prob(A\cap B\mid A\cup B)=1/3$. 因 $1/3\ne4/9$, 不独立. 原模型独立并不保证经过任意共同筛选后仍独立; 筛选已排除了二者都不发生的结果.
\end{solution}

\begin{exercise}
构造三个事件, 它们的三重交集概率等于三个概率的乘积, 但它们并不成对独立.
\end{exercise}
\begin{solution}
公平抛两枚硬币, 令 $A=B$ 为第一枚正面事件, $C=\varnothing$. 有 $\Prob(A\cap B\cap C)=0=\Prob(A)\Prob(B)\Prob(C)$, 但 $\Prob(A\cap B)=1/2\ne1/4$. 这表明只检查全交集的乘积公式不够. 若要避免零概率事件, 可在八个等可能结果 $\{1,\ldots,8\}$ 上取
\[
 A=\{1,2,3,4\},\quad B=\{1,2,3,5\},\quad C=\{1,6,7,8\}.
\]
三事件各概率 $1/2$, 三重交集只有元素一, 概率 $1/8$, 但 $A\cap B$ 有三个元素, 概率 $3/8\ne1/4$.
\end{solution}
